\documentclass[11pt]{article}

\usepackage[T1]{fontenc}
\usepackage{lmodern}
\usepackage{amsmath,amssymb,amsthm}
\usepackage{booktabs}
\usepackage[margin=1in]{geometry}
\usepackage[hidelinks]{hyperref}
\hypersetup{pdftitle={The Least Positive Layer of the P4--Paw Blow-Up Coefficient System},pdfauthor={Xiaojian Zeng},pdfsubject={Exact least-layer classification for a distance-cospectral blow-up coefficient system},pdfkeywords={distance spectrum, graph blow-up, exact enumeration, SMT, positive integer solutions}}
\usepackage{microtype}

\newtheorem{theorem}{Theorem}
\newtheorem{lemma}[theorem]{Lemma}
\newtheorem{proposition}[theorem]{Proposition}
\newcommand{\Zpos}{\mathbb Z_{>0}}
\newcommand{\Pfour}{P_4}

\title{The Least Positive Layer of the $P_4$--Paw
Blow-Up Coefficient System}
\author{Xiaojian Zeng\\
\small PBC School of Finance, Tsinghua University, Beijing 100190, China\\
\small \texttt{xiaojianzeng@foxmail.com}}
\date{8 September 2026}

\begin{document}
\maketitle

\begin{abstract}
Gong, Wang, and Wang exhibited a positive-integer solution of common sum
$61$ to the coefficient system that makes an independent-set blow-up of
$P_4$ distance-cospectral with a blow-up of the paw graph.  We determine the
first nonempty layer of that fixed system.  Its least common sum is $61$, and
the complete ordered layer consists of four solutions, forming one orbit
under the two system-preserving involutions.  The proof is computer-assisted
but exact: an exhaustive signature join covers all $521855$ positive ordered
quadruples on each side through sum $61$, and a separate bounded bit-vector
encoding confirms both the empty lower layers and the exhausted equality
layer.  This result concerns only the stated coefficient system; it neither
determines the least order among all connected distance-cospectral graph
pairs of opposite bipartiteness nor classifies higher solution layers.
\end{abstract}

\section{Context and scope}

For a connected graph, its distance spectrum is the spectrum of its distance
matrix.  Gong, Wang, and Wang used nonuniform independent-set blow-ups of
$\Pfour$ and the paw graph to construct connected distance-cospectral graphs
of opposite bipartiteness \cite{GongWangWang2026}.  Coefficientwise comparison
of their reduced distance characteristic polynomials gives the four equations
studied below.  Their displayed multiplicities yield one solution of common
sum $61$, but their paper does not exclude smaller positive solutions or
classify the solutions at sum $61$ \cite[Eq.~(3.3), Thm.~3.1 and
Conclusion]{GongWangWang2026}.

The distinction between an exact finite-domain theorem and its ambient
problem is important.  For comparison, exact finite-library results for
ternary Borsuk and kissing-code models are explicitly separated from the
corresponding unrestricted questions in \cite{ZengEtAl2026FiniteCertificates}.
Here the domain is likewise fixed: ordered positive-integer solutions of one
coefficient system, only through its first nonempty common-sum layer.  A
targeted search of primary arXiv records through 8 September 2026 found no
later version of \cite{GongWangWang2026} and no separate treatment of this
least layer.  This bounded search is not a claim of priority.

\section{Statement of the result}

Let $b=(b_0,b_1,b_2,b_3)$ and $a=(a_0,a_1,a_2,a_3)$ belong to
$\Zpos^4$.  The coefficient system is
\begin{align}
 b_0+b_1+b_2+b_3
   &=a_0+a_1+a_2+a_3,                                      \label{eq:sum}\\
 3b_0b_1-5b_0b_3+3b_1b_2+3b_2b_3
   &=3(a_0a_1+a_0a_2+a_0a_3+a_1a_2),                      \label{eq:quad}\\
 8b_0b_3(b_1+b_2)
   &=2a_1a_2(a_3-2a_0),                                    \label{eq:cubic}\\
 4b_0b_1b_2b_3
   &=a_0a_1a_2a_3.                                         \label{eq:quartic}
\end{align}
Write $s=\sum_i b_i=\sum_i a_i$.  Define the signatures
\begin{align*}
 \sigma_b(b)={}&\bigl(
  \textstyle\sum_i b_i,
  3b_0b_1-5b_0b_3+3b_1b_2+3b_2b_3,
  8b_0b_3(b_1+b_2),
  4b_0b_1b_2b_3\bigr),\\
 \sigma_a(a)={}&\bigl(
  \textstyle\sum_i a_i,
  3(a_0a_1+a_0a_2+a_0a_3+a_1a_2),
  2a_1a_2(a_3-2a_0),
  a_0a_1a_2a_3\bigr).
\end{align*}

\begin{theorem}\label{thm:main}
The least common sum of an ordered positive-integer solution of
\eqref{eq:sum}--\eqref{eq:quartic} is $61$.  At that sum the complete ordered
solution set is
\begin{equation}\label{eq:solutions}
 \begin{split}
 b&\in\{(1,16,8,36),(36,8,16,1)\},\\
 a&\in\{(4,9,16,32),(4,16,9,32)\},
 \end{split}
\end{equation}
with all four Cartesian combinations present.  Each has common signature
\[
 (61,1116,6912,18432).
\]
Moreover, the four ordered pairs form one orbit under precisely the two
allowed involutions
\[
 (b_0,b_1,b_2,b_3)\longmapsto(b_3,b_2,b_1,b_0),\qquad
 (a_0,a_1,a_2,a_3)\longmapsto(a_0,a_2,a_1,a_3).
\]
\end{theorem}

\section{Exact exhaustive proof}

\begin{lemma}[Signature equivalence]\label{lem:signature}
An ordered pair $(b,a)\in\Zpos^4\times\Zpos^4$ satisfies
\eqref{eq:sum}--\eqref{eq:quartic} if and only if
$\sigma_b(b)=\sigma_a(a)$.
\end{lemma}

\begin{proof}
Equality of the first signature coordinates is exactly \eqref{eq:sum};
equality of the next three coordinates is, in order,
\eqref{eq:quad}, \eqref{eq:cubic}, and \eqref{eq:quartic}.  Thus the
signature test is equivalent to the system, not a relaxation of it.
\end{proof}

For $s\geq4$, let
\[
 \mathcal C_s=\{x\in\Zpos^4:x_0+x_1+x_2+x_3=s\}.
\]
Choosing three cut points among $1,\ldots,s-1$ is a bijection with
$\mathcal C_s$, and hence
\begin{equation}\label{eq:counts}
 |\mathcal C_s|=\binom{s-1}{3},\qquad
 \sum_{s=4}^{61}|\mathcal C_s|=\binom{61}{4}=521855.
\end{equation}

\begin{proposition}[Exhaustive signature join]\label{prop:join}
Exact enumeration of $\mathcal C_s$ separately for the $a$- and $b$-side
signatures gives no equality $\sigma_b(b)=\sigma_a(a)$ for
$4\leq s\leq60$.  At $s=61$ it gives exactly the four pairs in
\eqref{eq:solutions}.
\end{proposition}

\begin{proof}[Certificate proof]
The accompanying program \texttt{enumerate\_signatures.py} generates a
quadruple by the nested ranges
\[
 \begin{gathered}
  1\leq x_0\leq s-3,\\
  1\leq x_1\leq s-x_0-2,\\
  1\leq x_2\leq s-x_0-x_1-1,\\
  x_3=s-x_0-x_1-x_2.
 \end{gathered}
\]
These ranges are another form of the cut-point bijection, so every member of
$\mathcal C_s$ occurs once and no other tuple occurs.  The program makes
separate passes over the two sides, evaluates the displayed signatures using
arbitrary-precision integers, and joins equal keys without pruning or
floating-point comparisons.

The report checks the count $\binom{s-1}{3}$ on every layer.  Its aggregate
results are summarized in Table~\ref{tab:join}.  The four reported pairs are
listed in Theorem~\ref{thm:main}; direct integer substitution gives their
common signature there.

\begin{table}[ht]
 \centering
 \caption{Exact signature-join coverage.}
 \label{tab:join}
 \begin{tabular}{@{}lrrr@{}}
  \toprule
  sums & $a$-tuples & $b$-tuples & ordered matches \\
  \midrule
  $4$--$60$ & $487635$ & $487635$ & $0$ \\
  $61$      & $34220$  & $34220$  & $4$ \\
  \midrule
  $4$--$61$ & $521855$ & $521855$ & $4$ \\
  \bottomrule
 \end{tabular}
\end{table}
\end{proof}

The join was checked by a structurally separate solver encoding.  This second
program does not import the tuple generator or its signature tables: it
asserts the eight coordinates and equations \eqref{eq:sum}--
\eqref{eq:quartic} directly in quantifier-free bit-vector logic and uses Z3
\cite{deMouraBjorner2008}.  Each coordinate is a six-bit unsigned integer in
$[1,58]$.  This cap is lossless, because four positive coordinates with sum
at most $61$ have each coordinate at most $58$.  Coordinates are zero-extended
to $32$ bits before arithmetic.

It remains to justify that modular arithmetic introduces no solutions.  On
$1\leq x_i\leq58$, the absolute differences between the two integer sides of
the sum, quadratic, cubic, and quartic equations are bounded respectively by
\[
 \begin{split}
  8\cdot58&=464<465,\\
  (14+12)58^2&=87464,\\
  (16+6)58^3&=4292464,\\
  (4+1)58^4&=56582480.
 \end{split}
\]
Here the coefficient sums come from taking absolute values term by term on
the two sides.  Every displayed bound is less than $2^{32}$.
Consequently, equality modulo $2^{32}$ is equivalent to ordinary integer
equality throughout the encoded domain, including the subtraction in
\eqref{eq:cubic}.

With common sum at most $60$, Z3 returned \texttt{unsat}.  With common sum at
most $61$, it returned the four pairs in \eqref{eq:solutions}; after a clause
blocking each complete ordered eight-tuple, the next result was
\texttt{unsat}.  Exact host-integer substitution checked every returned
model.  Thus the second route independently confirms both the empty lower
range and exhaustion after the fourth model.

\begin{proof}[Proof of Theorem~\ref{thm:main}]
Lemma~\ref{lem:signature} and Proposition~\ref{prop:join} exclude every
positive solution below $61$ and exhibit exactly four at $61$.  Hence $61$ is
the least common sum and \eqref{eq:solutions} is the complete ordered layer.

Reversing $b$ preserves every component of $\sigma_b$, and interchanging
$a_1,a_2$ preserves every component of $\sigma_a$.  The two involutions
commute.  Applied to the displayed representative
\[
 ((1,16,8,36),(4,9,16,32)),
\]
they produce the four distinct Cartesian combinations in
\eqref{eq:solutions}.  Completeness of that list therefore proves that the
layer is one permitted orbit.
\end{proof}

\section{Reproducibility and limitations}

The ancillary evidence consists of the exact enumeration program and its
JSON report, the independent bit-vector verifier and replayable SMT-LIB
instances, and a fail-closed comparison of the two model lists.  The recorded
run used CPython 3.10.12; the solver route used Z3 4.15.3 with random seed zero.
Neither route samples inputs, and no numerical tolerance occurs.  The scripts
\texttt{enumerate\_signatures.py}, \texttt{verify\_z3\_bv.py}, and
\texttt{compare\_verifiers.py} reproduce the three reports when run from the
project root in that order.

Theorem~\ref{thm:main} determines only the first positive layer of the fixed
$\Pfour$/paw coefficient system.  It does not classify solutions at sums
above $61$.  More importantly, exclusion within this ansatz is not exclusion
over all connected graphs: the least possible order of an arbitrary connected
distance-cospectral pair of opposite bipartiteness is not determined here.

\paragraph{Computational and AI assistance.}
OpenAI Codex assisted with targeted literature searching, implementation and
cross-check organization, and preparation of the manuscript.  Language-model
output is not used as mathematical evidence.  The theorem-bearing finite
claims are supported by the explicit domain reduction, exact programs,
machine-readable reports, and the independently implemented solver encoding.
The author takes responsibility for the statement and argument.

\bibliographystyle{abbrv}
\bibliography{references}

\end{document}
